\documentclass[12pt]{article}

\usepackage[a4paper, margin=1in]{geometry}
\usepackage{graphicx, parskip, setspace, ragged2e, enumitem}
\usepackage{amsmath, amssymb, amsthm, tikz, tikz-3dplot, chemfig}
\usepackage[english]{babel}
\setstretch{1.5}

\usepackage{accents}

\newcommand{\ut}[1]{\underaccent{\tilde}{#1}}
\renewcommand{\vec}[1]{\ut{#1}}

\newtheorem{theorem}{Theorem}[section]

\begin{document}
\center{\huge\textbf{Proof of the 3D version of the Euler Inequality}}

\vfill

\tdplotsetmaincoords{75}{0} % Set the view
\tdplotsetrotatedcoords{0}{0}{15}
\begin{tikzpicture}[tdplot_rotated_coords, scale=7]
	% Coordinates of tetrahedron vertices (normalized for the sphere)
	\coordinate (D) at ({sqrt(8/9)}, 0, -{1/3});
	\coordinate (C) at (-{sqrt(2/9)}, {sqrt(2/3)}, -{1/3});
	\coordinate (B) at (-{sqrt(2/9)}, -{sqrt(2/3)}, -{1/3});
	\coordinate (A) at (0, 0, 1);
	\coordinate (O) at (0, 0, 0);

	\filldraw[black] (0,0,0) circle (0.5pt);

	\draw[thick] (0,0,-{1/3}) circle ({2*sqrt(2)/3});
  \draw[thick] (0,0,-{sqrt(2)/6}) circle ({sqrt(2)/6});
  \tdplotsetrotatedcoords{80}{90}{0}
	\draw[thick, tdplot_rotated_coords] (0,0,0) circle (1);
  \draw[thick, tdplot_rotated_coords] (0,0,0) circle ({1/3});
	\tdplotsetrotatedcoords{0}{0}{15}

	% Draw the tetrahedron
	\draw[thick,black] (A) -- (B) -- (C) -- cycle; % Top face
	\draw[thick,black] (A) -- (C) -- (D) -- cycle; % Left face
	\draw[thick,black] (A) -- (D) -- (B) -- cycle; % Right face
	\draw[thick,black] (B) -- (D);          % Hidden edge

	% Label the tetrahedron vertices
	\node[above=5] at (A) {A};
	\node[below=5] at (B) {B};
	\node[above left] at (C) {C};
	\node[right=3, above=6] at (D) {D};
	\node[above right] at (O) {\textbf{O}};
	\filldraw [black] (A) circle (0.5pt);
	\filldraw [black] (B) circle (0.5pt);
	\filldraw [black] (C) circle (0.5pt);
	\filldraw [black] (D) circle (0.5pt);

  \draw[->, thick] (O) -- node[below] {\textbf{R}} (1,0,0);
  \draw[->, thick] (O) -- node[above left] {\textbf{r}} (-{1/3},0,0);

\end{tikzpicture}
\vskip 0.5cm
{\large{Diagram of the inscribed and circumscribed sphere of a tetrahedron}}

\vfill
\newpage

\justifying
\tableofcontents
\newpage

\section{Formal Statements of the Result}
\begin{theorem}[3D Euler Inequality]
  Let points $A$, $B$, $C$, $D$ denote the vertices of a tetrahedron in $\mathbb{R}^3$; let $R$, $r$ be the radii of the circumscribed and inscribed sphere respectively. Then,
  \[\forall A, B, C, D \in \mathbb{R}^3, R \geq 3r\text{.}\]
\end{theorem}
% end of formal Statements
\newpage

\section{Proof}
\begin{proof}[Proof of the inequality]
  We will first define the vector $\overrightarrow{OI}$, where $I$ is the center of the inscribed sphere, in terms of the vectors $\vec{a}$, $\vec{b}$, $\vec{c}$, $\vec{d}$ as the vectors to $A$, $B$, $C$, $D$ from the origin respectively.
\end{proof}
% end of proof

\end{document}
